\par\Needspace{12\baselineskip}\bigskip
\pdfbookmark[2]{Exposition, Section troisième}{pallas-II-section3}
\begin{center}
{\large\bfseries Section troisième.\par}
\medskip
\textit{Méthode de calculer les variations finies des élémens pendant un tems limité.}
\end{center}
\pdfbookmark[3]{Exposition, 17}{pallas-II-article17}
\begin{center}17.\end{center}
\noindent\hspace*{\parindent}La petitesse des masses perturbatrices comparées à la masse du soleil est une circonstance extrêmement favorable dans la théorie des perturbations planétaires. En toute rigueur, il faudrait, dans les calculs des variations \SourceLeaf{467}instantanées de chaque élément, employer les valeurs vraies des élémens qui entrent dans l’expression de ces variations. Cependant vu la petitesse de ces variations mêmes, le résultat sera fort peu différent, si l’on n’emploie que des valeurs approchées des élémens: du moins on obtient de cette manière une première approximation, qu’on pourra corriger dans la suite, si on le juge nécessaire.
\par Ainsi d’abord on déterminera, d’après les méthodes connues, des élémens purement elliptiques qui s’accordent le mieux possible avec les observations. Maintenant s’il ne s’agit que de la théorie du mouvement de la planète pendant un tems limité, on commencera par partager ce tems en plusieurs intervalles égaux, que nous supposerons chacun de \(\theta\) jours. Il sera convenable de choisir la première époque de ces intervalles, laquelle nous désignerons par \(T\), un peu avant le tems pendant lequel on veut calculer le mouvement, et de continuer la série \(T,\ T+\theta,\ T+2\theta,\ T+3\theta\) etc. un peu au delà de ce même tems. On calculera, d’après les formules développées ci-dessus, et sur les élémens approchés, les valeurs numériques des variations instantanées des différens élémens, pour les époques \(T,\ T+\theta,\ T+2\theta\) etc., savoir les valeurs de \(\dfrac{\mathrm{d}i}{\mathrm{d}t},\ \dfrac{\mathrm{d}\AscendingNode}{\mathrm{d}t},\ \dfrac{\mathrm{d}n}{\mathrm{d}t},\ \dfrac{\mathrm{d}\varphi}{\mathrm{d}t},\ \dfrac{\mathrm{d}\tilde\omega}{\mathrm{d}t},\ \dfrac{\mathrm{d}\epsilon}{\mathrm{d}t}\), car on n’a pas besoin de \(\dfrac{\mathrm{d}a}{\mathrm{d}t}\), la valeur de \(a\) étant toujours liée à celle de \(n\). Reste donc à intégrer ces six différentielles, et de déterminer, par une seconde intégration, le mouvement moyen \(\int n\mathrm{d}t\). Voici la méthode, qui nous paraît être la plus commode pour cet effet.
